\documentclass[lettersize,journal]{IEEEtran}

% ============================================================
% PACKAGES
% ============================================================
\usepackage{amsmath,amssymb,amsfonts,amsthm}
\usepackage{algorithm}
\usepackage{algpseudocode}
\usepackage{graphicx}
\usepackage{cite}
\usepackage{url}
\usepackage{xcolor}
\usepackage{booktabs}
\usepackage{multirow}
\usepackage{bm}
\usepackage[caption=false,font=footnotesize]{subfig}
\usepackage{balance}
\usepackage{tikz}
\usetikzlibrary{arrows.meta,positioning,calc,decorations.pathmorphing,
  shapes.geometric,fit,backgrounds,patterns,fadings,shadows.blur}
\usepackage{pgfplots}
\pgfplotsset{compat=1.18}

\hyphenation{op-tical net-works semi-conduc-tor Kolmo-go-rov}

% ============================================================
% THEOREM ENVIRONMENTS
% ============================================================
\newtheorem{corollary}{Corollary}
\newtheorem{lemma}{Lemma}
\newtheorem{remark}{Remark}
\newtheorem{definition}{Definition}
\newtheorem{proposition}{Proposition}
\newtheorem{assumption}{Assumption}

% ============================================================
% CUSTOM COMMANDS
% ============================================================
\DeclareMathOperator{\erfc}{erfc}
\DeclareMathOperator{\diag}{diag}
\DeclareMathOperator*{\argmin}{arg\,min}
\DeclareMathOperator*{\argmax}{arg\,max}
\newcommand{\vect}[1]{\mathbf{#1}}
\newcommand{\mat}[1]{\mathbf{#1}}
\newcommand{\set}[1]{\mathcal{#1}}
\newcommand{\norm}[1]{\left\|#1\right\|}
\newcommand{\abs}[1]{\left|#1\right|}
\newcommand{\R}{\mathbb{R}}
\newcommand{\E}{\mathbb{E}}

% ============================================================
\usepackage[active,tightpage]{preview}
\setlength\PreviewBorder{0.6pt}
\begin{document}
\begin{preview}
\begin{tikzpicture}
\begin{axis}[
  width=\columnwidth, height=5.4cm,
  ybar, bar width=13pt,
  every axis plot/.append style={bar shift=0pt},
  ylabel={$\epsilon_{\mathrm{recon}}$ (\%)},
  ymin=0, ymax=74, xmin=-0.7, xmax=5.7,
  xtick={0,1,2,3,4,5},
  xticklabels={linear, MM, Hill, hi-freq, bump, bimodal},
  x tick label style={font=\footnotesize, rotate=25, anchor=east},
  ytick={0,20,40,60},
  grid=major, grid style={gray!20},
  tick label style={font=\footnotesize}, label style={font=\normalsize},
  legend style={font=\footnotesize, at={(0.46,0.98)}, anchor=north, draw=gray!30, fill=white, fill opacity=0.9, text opacity=1},
]
\fill[red!8] (axis cs:-0.7,33) rectangle (axis cs:5.7,74);
\draw[red!50, densely dashed] (axis cs:-0.7,33) -- (axis cs:5.7,33);
\node[font=\footnotesize, text=red!60!black, anchor=north west] at (axis cs:-0.6,73) {reject region};
\addplot+[ybar, fill=blue!45, draw=blue!70, error bars/.cd, y dir=both, y explicit]
  coordinates {(0,3.9)+-(0,0.7) (1,7.0)+-(0,1.2) (2,29.1)+-(0,3.4)};
\addlegendentry{in-library}
\addplot+[ybar, fill=orange!60, draw=orange!80, error bars/.cd, y dir=both, y explicit]
  coordinates {(4,64.4)+-(0,0.3) (5,37.8)+-(0,1.6)};
\addlegendentry{out-of-lib.\ (resolved)}
\addplot+[ybar, fill=red!45, draw=red!70, error bars/.cd, y dir=both, y explicit]
  coordinates {(3,5.9)+-(0,1.8)};
\addlegendentry{out-of-lib.\ (sub-grid)}
\addplot[only marks, mark=*, mark size=0.9pt, draw=black!75, fill=black!75, forget plot]
  coordinates {(0,3.19) (0,4.83) (0,3.56) (1,5.91) (1,6.39) (1,8.77) (2,25.89) (2,27.68) (2,33.76) (3,7.33) (3,7.08) (3,3.31) (4,64.01) (4,64.66) (4,64.38) (5,39.5) (5,35.61) (5,38.17)};
\draw[->, gray!70, thick] (axis cs:3,22) -- (axis cs:3,9);
\node[font=\scriptsize, text=gray!70, align=center, anchor=south] at (axis cs:3,22) {smoothed:\\not flagged};
\end{axis}
\end{tikzpicture}
\end{preview}
\end{document}
